%%%%%%%%%%%%%%%%%%%%%%%%%%%%%%%%%%%%%%%
% This is a modified ONE COLUMN version of
% the following template:
%
% Deedy - One Page Two Column Resume
% LaTeX Template
% Version 1.1 (30/4/2014)
%
% Original author:
% Debarghya Das (http://debarghyadas.com)
%
% Original repository:
% https://github.com/deedydas/Deedy-Resume
%
% IMPORTANT: THIS TEMPLATE NEEDS TO BE COMPILED WITH XeLaTeX
%
% This template uses several fonts not included with Windows/Linux by
% default. If you get compilation errors saying a font is missing, find the line
% on which the font is used and either change it to a font included with your
% operating system or comment the line out to use the default font.
%%%%%%%%%%%%%%%%%%%%%%%%%%%%%%%%%%%%%%
\documentclass[]{deedy-resume-openfont}

\begin{document}
    
%%%%%%%%%%%%%%%%%%%%%%%%%%%%%%%%%%%%%%
%
%     Profile
%
%%%%%%%%%%%%%%%%%%%%%%%%%%%%%%%%%%%%%%
\namesection{Pingshi}{Yu}{pingshiyu@gmail.com | London, UK | linkedin.com/in/jacob-pingshi-yu/}

%%%%%%%%%%%%%%%%%%%%%%%%%%%%%%%%%%%%%%
%
%
%    Personal Statement
%
%%%%%%%%%%%%%%%%%%%%%%%%%%%%%%%%%%%%%%
%\section{Personal Statement}
%\raggedright
%\runsubsection{ }
%\begin{tightemize}
%	\item Fourth year studying Mathematics and Computer Science, with a strong academic background in both theoretical and applied areas of both subjects.
%	\item I am particularly interested in applications of machine learning to biology, and quantum computing algorithms. I am currently seeking PhD opportunities in a related field for further study.
%\end{tightemize}
%\sectionsep

%%%%%%%%%%%%%%%%%%%%%%%%%%%%%%%%%%%%%%
%
%     Education
%
%%%%%%%%%%%%%%%%%%%%%%%%%%%%%%%%%%%%%%
\section{Education}
\raggedright

\runsubsection{Imperial College London}\descript{| PhD in Computer Science}\hfill \location{Oct 2022 - Present}\\
Ongoing PhD in programming languages research, working on the modular validation of compilers, and property-based testing. \\
Topics of interest: algebraic effects, property-based testing, MLIR, complexity theory, functional programming\\
\sectionsep 

\runsubsection{University of Oxford}\descript{| Masters in Mathematics and Computer Science}\hfill \location{Oct 2019 - Jun 2020}\\
First class (80.3\%) \\
First class in all coursework, including: Quantum Computing, Probabilistic Algorithms, Computational Biology, Combinatorics. \\
First class dissertation in the intersection between graph theory, cheminformatics and machine learning. Title: "Systematic Pruning Methods for the Enumeration of Realistic Molecules".\\
\sectionsep 

\runsubsection{University of Oxford}\descript{| Undergrad in Mathematics and Computer Science}\hfill \location{Oct 2016 - Jul 2019}\\
Third Year - First Class; Second and First Year - Second Class, Division One\\
Awarded academic scholarship for high performance in Part B (third year) exams. \\
Broad range of courses from both Computer Science and Mathematics departments, including:\\ Computational Complexity, Graph Theory, Machine Learning, Logic and Set Theory Foundations, Continuous-Time Markov Chains
\sectionsep
  
\runsubsection{Dartford Grammar School}\descript{| IB and GCSE}\hfill \location{Sep 2009 - Jul 2016}\\
IB: 41/45 points (HL 777), HL Math, Physics and Computer Science. Subjects are marked out of 7.\\
\sectionsep

%%%%%%%%%%%%%%%%%%%%%%%%%%%%%%%%%%%%%%
%
%    Research Outputs
%
%%%%%%%%%%%%%%%%%%%%%%%%%%%%%%%%%%%%%%
\renewcommand\refname{PUBLICATIONS}
\nocite{*}
\bibliographystyle{unsrt}
\bibliography{citations}
\sectionsep

%%%%%%%%%%%%%%%%%%%%%%%%%%%%%%%%%%%%%%
%
%     Research Experience
%
%%%%%%%%%%%%%%%%%%%%%%%%%%%%%%%%%%%%%%
\section{Research Experience}
\runsubsection{Huawei Research Edinburgh}\descript{| Research Intern}\hfill \location{Edinburgh, UK | Oct 2025 - Current}
\begin{tightemize}
    \item Research into the implementation and application effect handlers for Huawei's open-source Cangjie programming language.
\end{tightemize}
\sectionsep

\runsubsection{Cambridge Quantum Computing}\descript{| Research Engineer}\hfill \location{Cambridge, UK | Jun - Aug 2020}
\begin{tightemize}
    \item Conducted systematic literature review on quantum computing using the ZX-calculus, a graphical calculus for reasoning about quantum systems.
    \item Gained a deep understanding, and completed successful implementation of scalable ZX-calculus (variant of ZX-calculus) in C++, for the representation, reasoning and simplification of large quantum circuits.
	\item Conducted systematical review of background literature, and successfully implemented a scalable graphical calculus for the representation and simplification of large quantum circuits.
	% \item Starting with no prior experience in C++: became proficient in C++ and familiarised with preexisting codebase within the first month.
	%\item 3000+ LoC\footnote{Lines of code, as counted by CLOC: http://cloc.sourceforge.net/} written with focus on readability, reusability and scalability, and is comprehensively documented.
\end{tightemize}
\sectionsep

\runsubsection{University of Oxford}\descript{| Masters Candidate}\hfill \location{Oxford, UK | Oct 2019 - Jul 2020}
\begin{tightemize}
    \item Masters research student under supervision Jotun Hein, Professor of Statistics and Bioinformatics.
    \item Conducted theoretical and numerical research into enumeration algorithms for the generation of synthetic chemical spaces.
    \item Implementation of novel techniques for the generation of chemical spaces, utilising orderly generation and machine learning on topological molecule descriptors.
    \item Work based on a chapter of thesis was later published at a top cheminformatics journal (10.1021/acs.jcim.0c01462)
\end{tightemize}
\sectionsep

\runsubsection{University of Oxford}\descript{| Research Assistant}\hfill \location{Oxford, UK | Apr - Jun 2018, Aug - Sep 2017}
\begin{tightemize}
	\item 2018: Invited back after a successful project previous summer. Designed questionnaires for data gathering, and for text sentimental algorithm validation.
	\item 2017: Research project in the Politics department investigating impact of legislations. Wrote a parser using Python and RegEx to process, filter and parse data from $\approx 10000$ statutory instruments. 
	\item Most valuable takeaway: research design, communication, technical documentation
\end{tightemize}
\sectionsep

%%%%%%%%%%%%%%%%%%%%%%%%%%%%%%%%%%%%%%
%
%     Talks/Summer Schools
%
%%%%%%%%%%%%%%%%%%%%%%%%%%%%%%%%%%%%%%
\section{Research Talks}
\raggedright
\runsubsection{\large{ASPLOS 2025}}\hfill \location{Rotterdam, NL | Apr 2025}\\
Gave a talk for our paper at ASPLOS 2025. Also presented at poster session. Website: https://www.asplos-conference.org/asplos2025/index.html/\\
\sectionsep

\runsubsection{\large{LaPRaS Seminar}}\hfill \location{London, UK | Dec 2024}\\
Talk on effect systems-based fuzzers in a departmental seminar. Website: https://lapras-imperial.github.io/\\
\sectionsep

\runsubsection{\large{ETH Zurich}}\hfill \location{Remote | Nov 2024}\\
Invited talk on composable fuzzers for programming languages using effect systems.\\
\sectionsep

\runsubsection{\large{VeTSS Summer School 2024}}\hfill \location{Bristol, UK | Aug 2024}\\
Overview of my research in the validation of composable compilers for the VeTSS Summer School. Website: https://wp.doc.ic.ac.uk/vss24/\\
\sectionsep

\runsubsection{\large{Imperial Computing Conference 2024}}\hfill \location{London, UK | Jul 2024}\\
Poster presentation on a validation framework for composable compilers and results. Won 3rd prize for poster presentation. \\
\sectionsep

\runsubsection{\large{S-REPLS 13}}\hfill \location{Bristol, UK | Nov 2023}\\
Talk about ongoing work during my PhD, on a framework to describe the semantics of a composable compiler IR (MLIR) and future directions. Website: https://plrg-bristol.github.io/fir/\\
\sectionsep

\runsubsection{\large{ISSTA 2023}}\hfill \location{Seattle, USA | Jul 2023}\\
Talk about the RustSmith work on fuzzing the Rust compiler. Gave both a presentation and a tool demonstration. Paper: https://dl.acm.org/doi/abs/10.1145/3597926.3604919 \\
\sectionsep

%%%%%%%%%%%%%%%%%%%%%%%%%%%%%%%%%%%%%%
%
%     Journal/Conference Reviews
%
%%%%%%%%%%%%%%%%%%%%%%%%%%%%%%%%%%%%%%
% \section{Journal / Conference Reviews}
% \raggedright
% \begin{tabular}{ l l }
% 	\descript{PL}   & {\location{FSE, FormaliSE}}
% \end{tabular}
% \sectionsep

%%%%%%%%%%%%%%%%%%%%%%%%%%%%%%%%%%%%%%
%
%     Skills
%
%%%%%%%%%%%%%%%%%%%%%%%%%%%%%%%%%%%%%%
\section{Skills}
\raggedright
\begin{tabular}{ l l }
	\descript{Programming}   & {\location{Skilled: Haskell, Python, C++; Familiar: Agda, Java, Scala, Perl, Matlab}} \\
	\descript{Technologies}  & {\location{Pandas, SQL, Unix, Tensorflow, \LaTeX, Git, HTML/CSS}} \\
	\descript{Proficiencies} & {\location{Data Science, Mathematics, Algorithms, Statistics, Machine learning}} \\
	\descript{Administrative} &{\location{College math dinners organiser, business consultant @ The Student Consultancy}} \\
	\descript{Languages} &{\location{Native proficiency in Mandarin Chinese and English}}
\end{tabular}
\sectionsep

%%%%%%%%%%%%%%%%%%%%%%%%%%%%%%%%%%%%%%
%
%	Work Experience
%
%%%%%%%%%%%%%%%%%%%%%%%%%%%%%%%%%%%%%%

\section{Work Experience}
\runsubsection{Credit Suisse}\descript{| Technical Analyst \& Qantitative Developer}\hfill \location{London, UK | Sep 2020 - Sep 2022}
\begin{tightemize}
    \item Quantitative developer for the market-making platform for foreign exchange.
    \item Lead the development of data pipelines for NLP-based market data into the pricing engine. Lead discussions with, and managed technical directions with quantitative, support and data teams.
	\item Creation various micro-services based tools for tasks including quantitative research, model validation, and release management.
    \item Conceptualisation, researching, and implementation of strategies to improve VM provisioning times through smart caching of VMs.
	\item Development of data driven pipelines and algorithms for monitoring \& optimizing code performance and developer productivity.
	\item Technologies and libraries used: Python, Docker, Jenkins, JavaScript, C++, Boost, Python, ElasticSearch
\end{tightemize}
\sectionsep
\runsubsection{Credit Suisse}\descript{| Technology Summer Analyst}\hfill \location{London, UK | Jun - Aug 2019}
\begin{tightemize}
    \item Technologist in the risk applications team, in the middle of the trading floor. Delivered 2 projects ahead of schedule in a fast paced environment. Both went into production and are in active use. 
	\item Discovered optimization in internal Perl API achieving $\approx 900 \times$ speedup in large SQL insert operations.
\end{tightemize}
\sectionsep
\runsubsection{ARM}\descript{| Data Engineering Intern}\hfill \location{Cambridge, UK | Jul - Sep 2018}
\begin{tightemize}
	\item Initiated, managed and delivered the backend (automated data engineering and processing) for a brand new company-wide capacity planning project.
	\item Most valuable takeaway: presentation skills, project management, data analytics, interpersonal skills
\end{tightemize}
\sectionsep

%%%%%%%%%%%%%%%%%%%%%%%%%%%%%%%%%%%%%%
%
%     Projects
%
%%%%%%%%%%%%%%%%%%%%%%%%%%%%%%%%%%%%%%
\section{Personal Projects}
\raggedright

\runsubsection{\large{Molecule Enumerator}}
\descript{| Java}\hfill \location{github.com/pingshiyu/openmg-multicore}\\
Extension of OMG\footnote{Open Molecule Generator, 10.1186/1758-2946-4-21}, with capabilities for enumerating chemical spaces and isomers under certain graph-theoretical constraints. Used in master's thesis and a subsequent publication. \\
\sectionsep

\runsubsection{\large{Greekletters}}
\descript{| Python}\hfill \location{github.com/pingshiyu/greekletters}\\
Convolutional neural network model trained on a Greek letters dataset, along with a GUI. Built to experiment with machine learning ideas, while working with a student funded start-up, in preparation for a Kaggle competition.\\
\sectionsep

%%%%%%%%%%%%%%%%%%%%%%%%%%%%%%%%%%%%%%
%
%     Hobbies
%
%%%%%%%%%%%%%%%%%%%%%%%%%%%%%%%%%%%%%%
\section{Hobbies}
\raggedright
\begin{tabular}{ l l }
	\descript{Chess}  & {Represented Oxford University in county \& varsity matches, ECF 174 (FIDE $\approx$2000)} \\
	\descript{Piano} & {Classically trained, grade 8. Working on Chopin \& movies/games music covers} \\
	\descript{Running} & {Ambitions of winning the local parkrun one day :) } \\
\end{tabular}
\sectionsep

\ 
\end{document}